\section{Threat Model}

In this section, we formalize the adversarial environment and trust boundaries within which CMatrix operates. Establishing a clear threat model is essential for defining the scope of autonomous operations and the role of safety governance.

\subsection{Adversary Profile and Capabilities}
We model the adversary as an AI-augmented red team or an autonomous security agent. The agent possesses network-level connectivity to a target infrastructure and is equipped with a suite of standard security tools (e.g., Nmap, Metasploit, SQLmap). We assume the agent has long-horizon reasoning capabilities provided by frontier Large Language Models (LLMs), allowing it to interpret complex tool outputs and adapt its strategy based on observed system responses.

\subsection{Adversary Goals and Assumptions}
The primary goal of the agent is the proactive identification and validation of technical vulnerabilities (e.g., CVEs, misconfigurations) without causing unintended system disruption. We assume:
\begin{itemize}
    \item The agent operates within a containerized "execution spine" with restricted network egress.
    \item The agent does not possess pre-existing knowledge of the target's internal topology.
    \item The target infrastructure may employ standard defensive measures such as IDS/IPS and Web Application Firewalls (WAF).
\end{itemize}

\subsection{Trust Boundary and Safety Governance}
As illustrated in Fig.~\ref{fig:threat-model}, we establish a formal \textbf{Trust Boundary} between the Reasoning Suite and the Target Infrastructure. This boundary is mediated by a risk-aware Human-in-the-Loop (HITL) safety gate.
\begin{itemize}
    \item \textbf{Internal Zone}: Contains the Orchestrator, ToT Strategy Selector, and ReWOO Planner. Operations in this zone are purely cognitive and do not interact with the target.
    \item \textbf{External Zone}: Contains the target network and the actual tool execution environment.
\end{itemize}
The safety gate enforces a risk-based policy $P$, where every proposed tool action $A$ is assigned a risk score $R_a(A)$. If $R_a(A)$ exceeds the vulnerability threshold, the orchestration is suspended, and human authorization is required before any network side-effects occur.

% ============================================================
% FIGURE 01: FORMAL THREAT MODEL
% File: figure-01-threat-model.tex
% Section: §3.1 Threat Model
% Description: Visual representation of the adversary, attack surface,
%              target system, and HITL safety boundaries.
% ============================================================

\begin{figure*}[htbp]
  \centering
  \begin{tikzpicture}[
    node distance=2cm,
    font=\small,
    box/.style={rectangle, draw, rounded corners, minimum width=2.5cm, minimum height=1cm, align=center, fill=gray!5},
    agent/.style={box, fill=blue!10, thick},
    gate/.style={diamond, draw, fill=red!10, minimum width=2.5cm, minimum height=1.5cm, align=center},
    target/.style={box, fill=green!10, thick},
    boundary/.style={draw, dashed, gray!60, inner sep=0.5cm}
  ]

    % Nodes
    \node (user) [box] {Security Operator \\ (Human-in-the-Loop)};
    \node (orch) [agent, below of=user, yshift=-0.5cm] {\textbf{CMatrix Orchestrator} \\ (Reasoning Suite)};
    \node (gate) [gate, right of=orch, xshift=2.5cm] {HITL \\ Safety Gate};
    \node (tools) [box, below of=gate, yshift=-0.5cm] {Security Tool Registry \\ (Nmap, Metasploit, etc.)};
    \node (target) [target, right of=gate, xshift=2.5cm] {Target Infrastructure \\ (Enterprise Network)};

    % Boundaries
    \begin{scope}[on background layer]
      \node [boundary, fit=(user) (orch) (gate) (tools), label=above:{\textbf{Trust Boundary}}] (trust) {};
    \end{scope}

    % Flows
    \draw[->, thick] (user) -- node[right] {Objective} (orch);
    \draw[->, thick] (orch) -- node[above, align=center] {Proposed \\ Actions} (gate);
    \draw[->, thick] (gate) -- node[left, align=center] {Approved \\ Tools} (tools);
    \draw[->, thick, red!60] (gate) -- node[above] {Veto} (user);
    \draw[->, thick] (tools) -- node[above, align=center] {Exploitation \\ Probes} (target);
    \draw[->, dashed] (target) -- node[below] {Observations} (orch);
    \draw[->, dashed] (orch) -- node[above right, align=center] {Risk Telemetry} (user);

    % Legend
    \node[anchor=north west] at (current bounding box.south west) [yshift=-0.5cm] {
      \begin{tabular}{ll}
        \tikz \draw[->, thick] (0,0) -- (0.5,0); & Control Flow \\
        \tikz \draw[->, dashed] (0,0) -- (0.5,0); & Data/Observation Flow \\
      \end{tabular}
    };

  \end{tikzpicture}
  \caption{Formal threat model of CMatrix. The security operator defines objectives for the autonomous orchestrator, while the HITL safety gate provides a hard trust boundary between the reasoning suite and the target infrastructure.}
  \label{fig:threat-model}
\end{figure*}

